\begin{proposition}[7.9.5]
\label{IV.2.s7.9.5-fr}
Sea $X$ un preesquema localmente noetheriano tal que, para todo preesquema $Y$
íntegro y finito sobre $X$, se pueda resolver $Y$. Entonces los anillos de los abiertos
afines de $X$ satisfacen las condiciones \textup{(ii)} y \textup{(iii)} de \textup{(7.8.2)}; si además
$X$ es universalmente catenario \textup{(5.6.3) (y}
\oldpage[IV-2]{220}
\textup{en particular si $X$ \emph{admite localmente una inmersión en un preesquema
regular} (5.6.4))}, $X$ es un preesquema excelente \textup{(7.8.5)}.
\end{proposition}

\begin{proof}
Esto se deduce de inmediato de (7.9.4) y (7.9.2).
\end{proof}
